\documentclass[a4paper,10pt]{article}
\usepackage[left=0.6in,top=0.5in,right=0.6in,bottom=0.5in]{geometry}
\usepackage{hyperref}
\usepackage{enumitem}
\usepackage{titlesec}
\usepackage{xcolor}
\usepackage{url}

% Define colors
\definecolor{darkgray}{RGB}{50,50,50}

% Formatting sections
\titleformat{\section}{\normalsize\bfseries\scshape}{}{0em}{}[\titlerule]
\titlespacing{\section}{0pt}{10pt}{6pt}

% Formatting lists
\setlist[itemize]{noitemsep, topsep=2pt, leftmargin=1.5em}

% Remove paragraph indentation
\setlength{\parindent}{0pt}

\begin{document}

\pagestyle{empty} % No page numbers

% Header
\begin{center}
    {\LARGE \textbf{Ayanda Thabo Phaketsi}} \\[4pt]
    \small
    \href{mailto:phaketsiayanda@gmail.com}{phaketsiayanda@gmail.com} $|$ 
    \href{https://phaketsi.vercel.app}{Portfolio: phaketsi.vercel.app} $|$ 
    \href{https://github.com/Past-da-king}{GitHub: Past-da-king} $|$ 
    \href{https://www.linkedin.com/in/ayanda-phakesti-747877248/}{LinkedIn} $|$ 
    South Africa
\end{center}

\vspace{4pt}

% Education
\section*{Education}

\textbf{University of Cape Town} \hfill Cape Town, South Africa \\
\textit{BSc Computer Science and Business Computing} \hfill \textbf{Dec 2025}

\vspace{4pt}
\textbf{Relevant Coursework:} Software Engineering, Algorithms and Data Structures, Databases, Operating Systems, Computer Networks, Machine Learning, Systems Analysis and Design

\vspace{8pt}

% Projects
\section*{Projects}
% Strat Edge Portal
\textbf{Strat Edge Portal Pro | Project Management System} \hfill \textbf{2026}\\
\small{Developed for Strat Edge}

\textit{Project Description:} 
Built a portal for tracking project schedules and finances. The system includes portfolio monitoring and reporting for stakeholders.

\textbf{Key Contributions:}
\begin{itemize}
    \item \textbf{System Build:} Created the platform using FastAPI and React with JWT authentication and access control.
    \item \textbf{Analytics:} Built backend services for project metrics, budget burndown, and task completion.
    \item \textbf{Storage:} Integrated document storage with Google Cloud Storage using signed URLs for access.
    \item \textbf{Dashboards:} Built charts and network diagrams to track project progress.
\end{itemize}

\vspace{8pt}
\textbf{Technologies Used:}
\begin{itemize}
    \item \textbf{Backend:} Python (FastAPI), SQLAlchemy, Google Cloud Storage, Pandas
    \item \textbf{Frontend:} React (TypeScript), Vite, Tailwind CSS, Recharts
\end{itemize}

\vspace{8pt}

% clia
\textbf{clia | Local AI Assistant for Software Tasks} \hfill \textbf{2025}\\
\small\url{https://pypi.org/project/clia-swe-ai/}

\textit{Project Description:} 
Built an AI assistant that helps with software tasks like refactoring and debugging by interacting with the local developer environment.

\textbf{Key Contributions:}
\begin{itemize}
    \item \textbf{Logic:} Built a loop for managing tool usage and multi-turn chat with Google Gemini.
    \item \textbf{Context:} Used Model Context Protocol (MCP) to build a tool server for filesystem and command-line actions.
    \item \textbf{Safety:} Added a permission model where users approve destructive operations.
    \item \textbf{Tooling:} Created tools for code editing and codebase snapshots.
\end{itemize}

\textbf{Technologies Used:}
\begin{itemize}
    \item \textbf{Languages/Frameworks:} Python, Google Gemini API, FastAPI, MCP
    \item \textbf{UI Tools:} Rich, prompt-toolkit, WebSockets
\end{itemize}

\vspace{8pt}

% MemGallery
\textbf{MemGallery | AI Memory Assistant for Android} \hfill \textbf{2025}\\
\small\url{https://github.com/Past-da-king/MemGallery}

\textit{Project Description:} 
Built an Android app that captures and processes media to enable semantic search, making AI features available on all Android 7.0+ devices.

\textbf{Key Contributions:}
\begin{itemize}
    \item \textbf{Capture:} Built a system to detect and process screenshots automatically.
    \item \textbf{Processing:} Integrated Gemini 2.5 Flash to extract events and notes from images and audio.
    \item \textbf{Search:} Built a memory bank where users find items using natural language.
    \item \textbf{Interface:} Built the UI using Material Design 3 and Jetpack Compose.
\end{itemize}

\textbf{Technologies Used:}
\begin{itemize}
    \item \textbf{Mobile Stack:} Kotlin, Jetpack Compose, Hilt, Room Database
    \item \textbf{System APIs:} WorkManager, ContentObservers, Google Generative AI SDK
\end{itemize}

\vspace{8pt}

% Nexa
\textbf{Nexa | Offline Communication Platform} \hfill \textbf{2025}\\
\small\url{https://github.com/Past-da-king/rednexa}

\textit{Project Description:} 
Worked with a team to build a communication system for areas without internet. The app allows messaging through a peer-to-peer network.

\textbf{Key Contributions:}
\begin{itemize}
    \item \textbf{Networking:} Built a custom protocol for data transfer via Bluetooth and Wi-Fi Direct.
    \item \textbf{Reliability:} Added logic for data deduplication and message delivery.
    \item \textbf{Security:} Integrated Google Tink for end-to-end encryption of messages.
    \item \textbf{Identity:} Created a local system for identity management without phone numbers.
\end{itemize}

\textbf{Technologies Used:}
\begin{itemize}
    \item \textbf{Core Stack:} Kotlin, Google Nearby Connections, Google Tink, Jetpack Compose
    \item \textbf{Data Management:} Room, Kotlinx Serialization
\end{itemize}

\vspace{8pt}

% StudyInc
\textbf{StudyInc | Collaborative Productivity Tool} \hfill \textbf{2025}\\
\small\url{https://phaketsi.vercel.app/#projects}

\textit{Project Description:} 
Built a collaborative platform for study sessions. It includes real-time presence, timers, and analytics to help students track their work.

\textbf{Key Contributions:}
\begin{itemize}
    \item \textbf{Syncing:} Used Socket.io for chat and timer sync, and WebRTC for video calls.
    \item \textbf{Desktop:} Used Electron to build features like overlays for timers and tasks.
    \item \textbf{Analytics:} Built a dashboard using React and Tailwind CSS to show study habits.
    \item \textbf{Backend:} Built a Node.js server with MongoDB and JWT authentication.
\end{itemize}

\textbf{Technologies Used:}
\begin{itemize}
    \item \textbf{Tech Stack:} Node.js, Next.js, Electron, MongoDB, Socket.io, WebRTC
    \item \textbf{Styling:} Tailwind CSS, shadcn/ui
\end{itemize}

\vspace{8pt}

% Verilookie
\textbf{Verilookie | Deepfake and AI Detection} \hfill \textbf{2025}\\
\small\url{https://github.com/Maphuti-Shilabje/verilookie}

\textit{Project Description:} 
Built a web app to detect deepfakes and AI-generated media, combined with a learning platform.

\textbf{Key Contributions:}
\begin{itemize}
    \item \textbf{Detection:} Integrated Hugging Face Vision Transformers and NVIDIA APIs for media analysis.
    \item \textbf{Education:} Built a quiz system powered by Gemini API that generates questions based on user history.
    \item \textbf{Optimization:} Used OpenCV for video frame sampling and Docker for containerization.
\end{itemize}

\textbf{Technologies Used:}
\begin{itemize}
    \item \textbf{Backend:} Python (FastAPI), PyTorch, OpenCV, Docker
    \item \textbf{Frontend/AI:} React, Google Gemini API, Hugging Face Transformers
\end{itemize}

\vspace{10pt}

% Awards & Achievements
\section*{Awards \& Achievements}

\begin{itemize}
    \item \textbf{Winner | Entelect University Challenge 2025:} Recognized for technical problem-solving.
    \item \textbf{2nd Place | Predictive Insights 2025 Hackathon:} Built a solution for data modeling.
\end{itemize}

\vspace{10pt}

% Technical Skills
\section*{Technical Skills}

\textbf{Languages:} Java, Python, JavaScript, TypeScript, Kotlin, C\#, SQL, \LaTeX \\
\textbf{Frameworks/Tools:} React, Next.js, FastAPI, Node.js, Jetpack Compose, Electron, MongoDB, Docker, Git \\
\textbf{Specializations:} AI/LLM Integration, Distributed Systems, Full-Stack Mobile \& Web

\end{document}
